\Needspace{6\baselineskip}\Article{345}
\textsc{Satz.} \textit{Es seien \((f,\allowbreak \lambda),\allowbreak (f,\allowbreak \mu)\) zwei gleichartige, identische oder verschiedene Perioden,} \textit{und \((f,\allowbreak \lambda)\) bestehe aus den Wurzeln \([\lambda],\allowbreak [\lambda'],\allowbreak [\lambda'']\) usw. Dann ist das Produkt aus \((f,\allowbreak \lambda)\) und \((f,\allowbreak \mu)\)} \textit{eine Summe von \(f\) gleichartigen Perioden, nämlich} \[=(f,\lambda+\mu)+(f,\lambda'+\mu)+(f,\lambda''+\mu)+\mathrm{usw.}=W\]

\textit{Beweis.} Es sei wie oben \(n-1=ef\); \(g\) sei eine primitive Wurzel für den Modul \(n\) und \(h=g^e\), woraus nach dem Vorhergehenden \((f,\allowbreak \lambda)=(f,\allowbreak \lambda h)=(f,\allowbreak \lambda hh)\) usw. folgt. Das gesuchte Produkt ist daher \[=[\mu].(f,\lambda)+[\mu h].(f,\lambda h)+[\mu hh].(f,\lambda hh)+\mathrm{usw.}\] und somit \[\begin{array}{@{}l@{\;}l@{\;}l@{}}
=[\lambda+\mu]&+[\lambda h+\mu]&\ldots+[\lambda h^{f-1}+\mu]\\
+[\lambda h+\mu h]&+[\lambda hh+\mu h]&\ldots+[\lambda h^f+\mu h]\\
+[\lambda hh+\mu hh]&+[\lambda h^3+\mu hh]&\ldots+[\lambda h^{f+1}+\mu hh]\ \mathrm{usw.}
\end{array}\] Dieser Ausdruck enthält insgesamt \(ff\) Wurzeln. Summiert man hier jede senkrechte Spalte für sich, so erhält man offenbar \[(f,\lambda+\mu)+(f,\lambda h+\mu)+\ldots+(f,\lambda h^{f-1}+\mu)\] Man sieht ohne Mühe, dass dieser Ausdruck mit \(W\) übereinstimmt, da die Zahlen \(\lambda,\allowbreak \lambda',\allowbreak \lambda''\) usw. nach Voraussetzung den Zahlen \(\lambda,\allowbreak \lambda h,\allowbreak \lambda hh\ldots\lambda h^{f-1}\) nach dem Modul \(n\) kongruent sein müssen (in welcher Reihenfolge, ist hier gleichgültig) und daher auch \[\begin{gathered}
\lambda+\mu,\ \lambda'+\mu,\ \lambda''+\mu\quad\mathrm{usw.}\quad\text{den Zahlen}\\
\lambda+\mu,\ \lambda h+\mu,\ \lambda hh+\mu\ldots\lambda h^{f-1}+\mu.\quad\textit{Q. E. D.}
\end{gathered}\]
\SourceLeaf{428}

Wir fügen diesem Satz die folgenden Folgerungen hinzu:

I. Bezeichnet \(k\) eine beliebige ganze Zahl, so ist das Produkt aus \((f,\allowbreak k\lambda)\) und \((f,\allowbreak k\mu)\) \[=(f,k(\lambda+\mu))+(f,k(\lambda^{\prime}+\mu))+(f,k(\lambda^{\prime\prime}+\mu))+\text{usw.}\]

II. Da die einzelnen Teile, aus denen \(W\) besteht, entweder mit der Summe \((f,\allowbreak 0)\), die \(=f\) ist, oder mit einer von \((f,\allowbreak 1),\allowbreak \ (f,\allowbreak g),\allowbreak \ (f,\allowbreak gg)\ldots(f,\allowbreak g^{e-1})\) übereinstimmen, lässt sich \(W\) auf die folgende Form bringen \[\begin{aligned}
W=af+b(f,1)+b^{\prime}(f,g)&+b^{\prime\prime}(f,gg)\\
&+\ldots+b^\varepsilon(f,g^{e-1})
\end{aligned}\] Dabei sind die Koeffizienten \(a,\allowbreak b,\allowbreak b^{\prime}\) usw. positive ganze Zahlen (oder einige auch \(=0\)). Ferner ist klar: Das Produkt aus \((f,\allowbreak k\lambda)\) und \((f,\allowbreak k\mu)\) wird dann zu \[=af+b(f,k)+b^{\prime}(f,kg)+\ldots+b^\varepsilon(f,kg^{e-1})\] So wird \textit{zum Beispiel} für \(n=19\) das Produkt aus der Summe \((6.1)\) und sich selbst oder das Quadrat dieser Summe \(=(6,\allowbreak 2)+(6,\allowbreak 8)+(6,\allowbreak 9)+(6.12)+(6,\allowbreak 13)+(6,\allowbreak 19)=6+2(6,\allowbreak 1)+(6,\allowbreak 2)+2(6,\allowbreak 4).\)

III. Da sich das Produkt aus jedem einzelnen Teil von \(W\) und einer gleichartigen Periode \((f,\allowbreak \nu)\) auf eine entsprechende Form bringen lässt, ist klar, dass auch das Produkt aus den drei Perioden \((f,\allowbreak \lambda).(f,\allowbreak \mu).(f,\allowbreak \nu)\) durch \(cf+d(f,\allowbreak 1)\ldots+d^\varepsilon(f,\allowbreak g^{e-1})\) dargestellt werden kann und die Koeffizienten \(c,\allowbreak d\) usw. ganze und positive Zahlen (oder \(=0\)) sind; außerdem gilt für jeden beliebigen ganzzahligen Wert von \(k\) \[(f,k\lambda).(f,k\mu).(f,k\nu)=cf+d(f,k)+d^{\prime}(f,kg)+\text{usw.}\] Ebenso erstreckt sich dieser Satz auf Produkte aus beliebig vielen gleichartigen Perioden, und es ist gleichgültig, ob diese Perioden sämtlich verschieden oder zum Teil oder sämtlich identisch sind.

IV. Daraus folgt: Werden in irgendeiner ganzrationalen algebraischen Funktion \(F=\varphi(t,\allowbreak u,\allowbreak v\ldots)\) für die Unbestimmten \(t,\allowbreak u,\allowbreak v\) usw. jeweils die gleichartigen Perioden \((f,\allowbreak \lambda),\allowbreak (f,\allowbreak \mu),\allowbreak (f,\allowbreak \nu)\) usw. eingesetzt, so lässt sich ihr Wert auf die Form \[A+B(f,1)+B^{\prime}(f,g)+B^{\prime\prime}(f,gg)+\ldots+B^\varepsilon(f,g^{e-1})\] bringen, und die Koeffizienten \(A,\allowbreak B,\allowbreak B^{\prime}\) usw. sind sämtlich ganze Zahlen, wenn alle festen Koeffizienten in \(F\) ganze Zahlen sind. Werden danach für \(t,\allowbreak u,\allowbreak v\) usw. jeweils
\SourceLeaf{429}
\((f,\allowbreak k\lambda),\allowbreak (f,\allowbreak k\mu),\allowbreak (f,\allowbreak k\nu)\) usw. eingesetzt, so führt sich der Wert von \(F\) auf \(A+B(f,\allowbreak k)+B^{\prime}(f,\allowbreak kg)+\text{usw.}\) zurück.

\Needspace{6\baselineskip}\Article{346}
\textsc{Satz.} \textit{Nimmt man an, dass \(\lambda\) eine durch \(n\) nicht teilbare Zahl ist, und schreibt man} \textit{der Kürze wegen \(p\) für \((f,\allowbreak \lambda)\), so lässt sich jede andere gleichartige Periode \((f,\allowbreak \mu)\), wobei auch \(\mu\) als durch \(n\)} \textit{nicht teilbar vorausgesetzt wird, auf die folgende Form bringen} \[\alpha+\beta p+\gamma pp+\ldots+\theta p^{e-1}\] \textit{wobei die Koeffizienten \(\alpha,\allowbreak \beta\) usw. feste rationale Größen sind.}

\textit{Beweis.} Wir bezeichnen zur Abkürzung die Perioden \((f,\allowbreak \lambda g),\allowbreak (f,\allowbreak \lambda gg),\allowbreak (f,\allowbreak \lambda g^3)\) usw. bis \((f,\allowbreak \lambda g^{e-1})\), deren Anzahl \(e-1\) ist und mit deren einer \((f,\allowbreak \mu)\) notwendig übereinstimmt, mit \(p^{\prime},\allowbreak p^{\prime\prime},\allowbreak p^{\prime\prime\prime}\) usw. Man hat dann unmittelbar die Gleichung \[0=1+p+p^{\prime}+p^{\prime\prime}+p^{\prime\prime\prime}+\text{usw.}\quad\mathinner{\ldotp\ \ldotp\ \ldotp\ \ldotp\ \ldotp}\quad(\mathrm I)\] Entwickelt man aber nach den Regeln des vorigen Artikels die Werte der Potenzen von \(p\) bis zur Potenz mit dem Exponenten \(e-1\), so ergeben sich \(e-2\) weitere Gleichungen der folgenden Art. \[\begin{aligned}0&=pp+A+ap+a^{\prime} p^{\prime}+a^{\prime\prime} p^{\prime\prime}+a^{\prime\prime\prime} p^{\prime\prime\prime}+\text{usw.}\quad\mathinner{\ldotp\ \ldotp}\quad(\mathrm{II})\\0&=p^3+B+bp+b^{\prime} p^{\prime}+b^{\prime\prime} p^{\prime\prime}+b^{\prime\prime\prime} p^{\prime\prime\prime}+\text{usw.}\quad\mathinner{\ldotp\ \ldotp}\quad(\mathrm{III})\\0&=p^4+C+cp+c^{\prime} p^{\prime}+c^{\prime\prime} p^{\prime\prime}+c^{\prime\prime\prime} p^{\prime\prime\prime}+\text{usw.}\quad\mathinner{\ldotp\ \ldotp}\quad(\mathrm{IV})\quad\text{usw.}\end{aligned}\] Dabei sind sämtliche Koeffizienten \(A,\allowbreak a,\allowbreak a^{\prime}\) usw., \(B,\allowbreak b,\allowbreak b^{\prime}\) usw. usw. ganze Zahlen und, was sorgfältig zu beachten ist und aus dem vorigen Artikel unmittelbar folgt, von \(\lambda\) völlig unabhängig; \textit{das heißt}, dieselben Gleichungen gelten auch dann noch, welchen anderen Wert man \(\lambda\) auch gibt. Diese Bemerkung erstreckt sich offenbar auch auf Gleichung I, sofern man \(\lambda\) als durch \(n\) nicht teilbar annimmt. --- Nehmen wir \((f,\allowbreak \mu)=p^{\prime}\) an; denn man sieht sehr leicht, dass, wenn \((f,\allowbreak \mu)\) mit einer anderen der Perioden \(p^{\prime\prime},\allowbreak p^{\prime\prime\prime}\) usw. übereinstimmt, zu den folgenden völlig analoge Schlüsse angewendet werden können. Da die Anzahl der Gleichungen I, II, III usw. \(e-1\) ist, lassen sich die Größen \(p^{\prime\prime},\allowbreak p^{\prime\prime\prime}\) usw., deren Anzahl \(=e-2\) ist, nach bekannten Methoden daraus eliminieren, sodass sich eine von ihnen freie Gleichung der folgenden Art \((Z)\) ergibt \[0=\mathfrak A+\mathfrak Bp+\mathfrak Cpp+\text{usw.}+\mathfrak Mp^{e-1}+\mathfrak Np^{\prime}\] Dies kann so geschehen, dass alle Koeffizienten \(\mathfrak A,\allowbreak \mathfrak B\ldots\mathfrak N\) ganze Zahlen und sicher nicht alle \(=0\) sind. Gilt hier nun nicht \(\mathfrak N=0\), so ist unmittelbar klar, dass sich \(p^{\prime}\) daraus in der im
\SourceLeaf{430}
Satz angegebenen Weise bestimmt. Es bleibt also zu zeigen, dass \(\mathfrak N=0\) nicht eintreten kann.

Nimmt man \(\mathfrak N=0\) an, so wird Gleichung \(Z\) zu \(\mathfrak Mp^{e-1}+\text{usw.}+\mathfrak Bp+\mathfrak A=0\); da sie sicher nicht über den Grad \(e-1\) hinausgeht, können ihr nicht mehr als \(e-1\) verschiedene Werte von \(p\) genügen. Da aber die Gleichungen, aus denen \(Z\) abgeleitet wurde, von \(\lambda\) unabhängig sind, ist klar, dass auch \(Z\) nicht von \(\lambda\) abhängt, also gilt, welche durch \(n\) nicht teilbare ganze Zahl man auch für \(\lambda\) nimmt. Daher wird Gleichung \(Z\) erfüllt, welcher von den \(e\) Summen \((f,\allowbreak 1),\allowbreak (f,\allowbreak g),\allowbreak (f,\allowbreak gg)\ldots(f,\allowbreak g^{e-1})\) man \(p\) auch gleichsetzt. Daraus folgt unmittelbar, dass diese Summen nicht sämtlich ungleich sein können, sondern wenigstens zwei einander gleich sein müssen. Die eine von zwei solchen gleichen Summen enthalte die Wurzeln \([\zeta],\allowbreak [\zeta^{\prime}],\allowbreak [\zeta^{\prime\prime}]\) usw., die andere \([\eta],\allowbreak [\eta^{\prime}],\allowbreak [\eta^{\prime\prime}]\) usw.; wir nehmen an (was zulässig ist), dass alle Zahlen \(\zeta,\allowbreak \zeta^{\prime},\allowbreak \zeta^{\prime\prime}\) usw., \(\eta,\allowbreak \eta^{\prime},\allowbreak \eta^{\prime\prime}\) usw. positiv und \(<n\) sind; offenbar sind sie auch sämtlich verschieden und keine \(=0\). Die Funktion \[x^\zeta+x^{\zeta^{\prime}}+x^{\zeta^{\prime\prime}}+\text{usw.}-x^\eta-x^{\eta^{\prime}}-x^{\eta^{\prime\prime}}-\text{usw.}\] hat kein höheres Glied als \(x^{n-1}\). Bezeichnen wir sie mit \(Y\), so gilt offenbar \(Y=0\), wenn man \(x=[1]\) setzt. Daher enthält \(Y\) den Faktor \(x-[1]\), den sie mit der im Vorhergehenden durch \(X\) bezeichneten Funktion \textit{gemeinsam} hat; dass dies aber absurd ist, lässt sich leicht zeigen. Hätte nämlich \(Y\) mit \(X\) irgendeinen gemeinsamen Faktor, so hätte der \textit{größte} gemeinsame Teiler der Funktionen \(X,\allowbreak Y\) (dass dieser sicher nicht bis zum Grad \(n-1\) reichen kann, geht schon daraus hervor, dass \(Y\) durch \(x\) teilbar ist) ausschließlich rationale Koeffizienten, wie aus der Art der Rechenoperationen zur Bestimmung des größten gemeinsamen Teilers zweier solcher Funktionen, deren Koeffizienten alle rational sind, unmittelbar folgt. In Artikel 341 haben wir aber gezeigt, dass \(X\) keinen Faktor eines Grades unter \(n-1\) enthalten kann, dessen Koeffizienten sämtlich rational sind. Daher kann die Annahme, dass \(\mathfrak N=0\) gilt, nicht bestehen.

\textit{Beispiel.} Für \(n=19,\allowbreak \ f=6\) gilt \(pp=6+2p+p^{\prime}+2p^{\prime\prime}\); daraus und aus \(0=1+p+p^{\prime}+p^{\prime\prime}\) folgt \(p^{\prime}=4-pp,\allowbreak \ p^{\prime\prime}=-5-p+pp\). Daher \[\begin{aligned}(6,2)&=4-(6,1)^2,\ &(6,4)&=-5-(6,1)+(6,1)^2\\(6,\allowbreak 4)&=4-(6,\allowbreak 2)^2,\allowbreak \ &(6.1)&=-5-(6,\allowbreak 2)+(6,\allowbreak 2)^2\\(6,\allowbreak 1)&=4-(6,\allowbreak 4)^2,\allowbreak \ &(6,\allowbreak 2)&=-5-(6,\allowbreak 4)+(6,\allowbreak 4)^2\end{aligned}\]
\SourceLeaf{431}


